\documentclass[11pt,a4paper]{article}
\usepackage[utf8]{inputenc}
\usepackage[T1]{fontenc}
\usepackage{hyperref}

% Manual margins instead of the geometry package (no external package needed)
\setlength{\topmargin}{-1.5cm}
\setlength{\oddsidemargin}{0cm}
\setlength{\evensidemargin}{0cm}
\setlength{\textwidth}{16cm}
\setlength{\textheight}{24cm}
\setlength{\headheight}{0cm}
\setlength{\headsep}{0cm}

\title{\textbf{Inteligentný systém na automatizované testovanie bezpečnosti webových aplikácií pomocou veľkých jazykových modelov}}
\author{}
\date{}



\begin{document}
\maketitle
\vspace{-1cm}

\noindent\textit{Využitie veľkých jazykových modelov a AI agentov na automatizované bezpečnostné testovanie webových aplikácií.}

\bigskip


\begin{center}
{\LARGE \textbf{AI-PENTEST}}\\[3mm]
Kyrylo Bukhanets, Artur Toba, Maksym Tomkevych , Andrii Tsybulskyi\\
FIIT STU\\
bukhanets@stuba.sk, xtoba@stuba.sk , xtomkevych@stuba.sk , xtribulskyi@stuba.sk\\
4.\ 10.\ 2026
\end{center}

\section*{Smer (oblasť výskumu)}

\noindent Aplikovaný výskum v oblasti umelej inteligencie a kybernetickej bezpečnosti.

\section*{Čo presne budeme robiť}

\noindent
Cieľom nie je iba to, že ,,neurónová sieť hľadá hackerské zraniteľnosti``. Vytvárame \textbf{Web Pentest Harness} -- špecializované prostredie, v ktorom AI agent dostane webovú aplikáciu a pokúsi sa samostatne vykonať jej kontrolované bezpečnostné testovanie.

\medskip
\noindent\textbf{Princíp:}

\begin{center}
\textit{Webová aplikácia $\rightarrow$ AI agent $\rightarrow$ Pentest Harness $\rightarrow$ testy $\rightarrow$ zistené zraniteľnosti $\rightarrow$ automatické vyhodnotenie}
\end{center}

\noindent\textbf{Harness bude:}
\begin{itemize}
\item poskytovať AI agentovi bezpečné testovacie prostredie;
\item poskytovať mu nástroje na interakciu s webovou aplikáciou;
\item zaznamenávať všetky akcie modelu;
\item overovať, či nájdená zraniteľnosť skutočne existuje;
\item hodnotiť úspešnosť agenta;
\item zbierať výsledky na následné porovnanie rôznych AI modelov.
\end{itemize}

\section*{Prečo je to dobrá téma}

\noindent
Ide o konkrétny inžiniersky problém: súčasné jazykové modely dokážu vykonávať zložité úlohy súvisiace s analýzou webových aplikácií, no chýba vhodné štandardizované prostredie na bezpečné a objektívne porovnávanie ich schopností v oblasti penetračného testovania. Navrhovaným riešením je práve vytvorenie Web Pentest Harness, ktorý umožní štandardizovane testovať a hodnotiť AI agentov na kontrolovaných webových zraniteľnostiach -- ide teda o skutočný aplikovaný výskum, nie len o ,,neurónovú sieť na hľadanie hackerských dier``.

\section*{Anotácia}

\noindent
Súčasné veľké jazykové modely dokážu vykonávať zložité úlohy súvisiace s analýzou webových aplikácií, no v súčasnosti chýba vhodné štandardizované prostredie na bezpečné a objektívne porovnávanie ich schopností v oblasti penetračného testovania. Projekt AI-PENTEST preto navrhuje vytvoriť Web Pentest Harness -- kontrolované prostredie, v ktorom AI agent dostane webovú aplikáciu a pokúsi sa samostatne vykonať jej bezpečnostné testovanie. Harness poskytuje agentovi bezpečné testovacie prostredie aj nástroje na interakciu s aplikáciou, zaznamenáva všetky jeho akcie, overuje, či nájdená zraniteľnosť skutočne existuje (prostredníctvom kontrolovaného proof-of-concept overenia), vyhodnocuje úspešnosť agenta a zbiera výsledky na následné porovnanie viacerých AI modelov. Jadrom riešenia je kontajnerizovaná sada zámerne zraniteľných webových aplikácií so známymi triedami chýb (SQL injection, XSS, autentifikačné a autorizačné obchádzky, SSRF, IDOR, nesprávna konfigurácia a pod.), doplnená anotovanou databázou referenčných zraniteľností (ground truth) a automatizovaným hodnotiacim modulom počítajúcim metriky ako presnosť, úplnosť (recall), mieru falošne pozitívnych nálezov a čas potrebný na nájdenie a overenie zraniteľnosti. Agent pracuje výlučne v sieťovo izolovanom prostredí s riadenými povoleniami, čo umožňuje bezpečné a reprodukovateľné experimenty. Výstupom projektu bude otvorená infraštruktúra (harness + benchmark sada úloh rôznej obtiažnosti) a empirické porovnanie viacerých architektúr AI agentov z hľadiska efektivity, spoľahlivosti a typických zlyhaní. Projekt tak prispieva k zodpovednému a štandardizovanému hodnoteniu kybernetických schopností AI systémov.

\section*{Právne zdroje}

\begin{enumerate}

\item OWASP Foundation. \textit{Web Security Testing Guide v 4.2}. 2020.\\
\url{https://owasp.org/www-project-web-security-testing-guide/}

\item OWASP Foundation. \textit{The OWASP Web Security Testing Framework}. 2020.\\
\url{https://wstg.owasp.org/stable/3-The_OWASP_Testing_Framework/0-The_Web_Security_Testing_Framework/}

\item OWASP Foundation. \textit{OWASP WebGoat}.\\
\url{https://owasp.org/www-project-webgoat/}

\item National Institute of Standards and Technology. \textit{AI Measurement and Evaluation}. 2024.\\
\url{https://www.nist.gov/ai-measurement-and-evaluation}

\item National Institute of Standards and Technology. \textit{Building Evaluation Probes into Agentic AI}. 2026.\\
\url{https://www.nist.gov/programs-projects/building-evaluation-probes-agentic-ai}

\item Jensen, T., Przybocki, M., Kodwani, M., Greene, K., Hall, P., Amironesei, R., Greenberg, C. \textit{ARIA Evaluation Planning Manual: Elements of ARIA-Style AI Evaluations}. NIST AI 200-3. 2026.\\
\url{https://doi.org/10.6028/NIST.AI.200-3}

\end{enumerate}

\end{document}

\end{document}
